% ISM 0.1
% Dense core 1
% both models larger grains 
% compare with other disks 

The maximum grain size for also varies with filling factor. As the filling factor decreases, the best-fit maximum grain size increases, ranging from $162$ - 535 \micron for the spherical models. This trend is likely due to the lower amount of material in a porous grain. As the filling factor decreases, a larger radius may be required to produce the same scattering properties of compact grains. Irregular grain models, where porosity is not varies, therefore, has a much smaller range of best-fit grain sizes across the fitting cases, ranging from 209-212\micronNS. The best-fit maximum dust grain size is \sbest for spherical grains with $f=0.5$, and \ibest for irregular grains.

 
Figure \ref{fig: Kataoka2015Figure3Recreation} showed that the \Pweff curves for compact grains do not peak at the characteristic grain size, but rather at a slightly larger grain size. Despite this slight difference due to varying filling factor, it can still be used to estimate the grain size the observations are most sensitive to. The characteristic grain sizes that correspond to the observed wavelength range from $\sim140-340$\micron (Table \ref{tab: amax_peak_values}). The best-fit grain sizes and many of the other fitted grain sizes, fall within this range of grain sizes expected to be detected at the observed wavelengths. 

% 156.67, 115. 67, 101.58, 69 


% Al all shape, porosity, and filling combinations the maximum fit grain size ranges from 162-535 \micronNS.  The characteristic grain sizes for the observed wavelengths range from $\sim140-340$ \micronNS, indicating that the best-fit grain size, and majority of the other fitted grain sizes, are consistent within the grain sizes we expected to be able to probe. 





